\section{Integrating the workers into Forge and Leant}
\label{sec:integration}

\subsection{The missing component is a semantic compiler}

The mathematical certificate checkers are intentionally small. The more subtle
Lean work is showing that the reified problem denotes the original goal. A
correct polynomial identity about an incorrectly translated expression proves
nothing useful. Forge's own contract already separates certificate soundness
from this correspondence~\cite{forge-contracts}; the new contribution is to
make the required correspondence concrete for these workers.

For a source state type $S$, input type $A$, and original update
$\operatorname{step}:S\to A\to S$, a positive bridge should contain:
\begin{align*}
 e &: S\to\Q^D,\\
 \iota &: A\to\Sigma,\\
 e(\operatorname{step}(s,a))&=F_{\iota(a)}(e(s)),\\
 e(\operatorname{initial}(\theta))&=S_0(\theta),\\
 q(e(s))=0&\Longrightarrow \operatorname{Goal}(s).
\end{align*}
Each displayed equality or implication is a Lean theorem, not a callback
assertion. These facts, the certificate theorem, and induction on the original
input list produce the source goal. Injectivity of $e$ is not required when the
observation implication already gives exactly the desired proposition.

A negative bridge needs more. A rational state or action produced by a solver
might have no source inhabitant. Thus a counterexample becomes a user-visible
refutation only after constructing the original typed parameter values and
inputs and checking the original proposition at that concrete execution. A
positive simulation into an overapproximation is insufficient to justify
lifting a negative result back.

\subsection{An admission algorithm for numeric folds}

The first reifier should target explicit folds over lists and a small first-order
state grammar. It should not begin by trying to translate arbitrary Lean
metaprograms or dependent recursive definitions. A practical admission pass is:

\begin{enumerate}
\item Normalize only the approved definitional equations of the fold and its
step function. Preserve lets until their dependencies and sharing are recorded.
\item Collect scalar state fields and parameters. Admit rational ring operations
and natural exponents. Translate each expression by an interpretation theorem,
so addition, multiplication, constants, and projection all come with source
correspondence.
\item Classify state dependence. Degree at most one yields Worker A after
homogenization; general admitted polynomial dependence yields Worker B.
\item Enumerate an actual finite input type, or use a supplied exhaustive
classification theorem. Merely seeing a few constructors in samples does not
establish exhaustive action coverage.
\item Build the target observation and initial parameterization. Preserve casts,
state-field order, and every parameter dependency explicitly.
\item Submit the exact reified problem. Replay the returned algebraic identities
and compose the semantic bridge with the source fold-induction theorem.
\end{enumerate}

Every rejected construct should be named: for example, ``state-dependent
integer division'', ``unknown numeric dictionary'', or ``input alphabet not
proved exhaustive''. Falling back to Forge's existing engines is legitimate;
pretending the unsupported expression is polynomial is not.

A source fold may have a nonnumeric payload alongside a numeric accumulator.
It can be erased only when a theorem establishes that the payload does not
affect the observation or the numeric update. If it controls a branch or carries
a typeclass method, it cannot be ignored. This is where a seemingly convenient
``just extract the arithmetic'' implementation often becomes unsound.

\subsection{Universal behavior as a new Leant verification branch}

Leant already checks an exact candidate against its type and then tries to
prove or refute a supplied behavioral proposition using \code{decide} and
bounded \code{simp}. Its documentation carefully distinguishes a conjunction
of examples from a universal assertion, and reports inconclusive checks rather
than inventing a counterexample~\cite{leant-behavior}. The proposed addition is
a restricted universal-behavior branch, not a replacement for that discipline.

For a candidate implementing a numeric fold and a reference specification, the
adapter should compile both to the admitted state model, form the product
observation, and ask Worker A or B to prove it zero for all input words. On
success, it reconstructs a theorem about the \emph{same exact candidate}. On
failure, a concrete word can become a new rejecting behavioral example after
source-level replay. On \Unknown, the candidate remains inconclusive for the
universal assertion.

This can improve synthesis in two different ways. A universal certificate can
accept a candidate that finite executable tests cannot establish correct.
Conversely, a distinguishing word can reject many other candidates cheaply,
provided each candidate is separately evaluated at that word. The latter is
ordinary counterexample-guided filtering; it does not transfer one candidate's
proof or rejection receipt to another.

The adapter should retain the already verified source identity, lexical
universe choices, and dictionaries. Djex and Leant invest substantial effort in
preserving precisely those associations~\cite{djex,leant}. This article does not
propose to rebuild their existing ownership or rank-N search mechanisms. It
adds a mathematical behavior verifier behind that boundary.

\subsection{A warning about Church encodings and parametricity}

A proof about folds of genuine finite lists does not automatically prove a
statement about every inhabitant of a Church-encoded type. In Lean, one cannot
silently import a Haskell-style free theorem, erase a universe distinction, or
assume that every polymorphic value is the image of an ordinary list encoding.
Those claims need their own source-level theorems and hypotheses.

The first integration should therefore accept either an explicit \code{List}
fold or a Church value accompanied by a proved representation theorem. A
certificate for all encoded finite lists should be displayed at that scope,
not as an answer to a more general original quantified signature. In particular,
this proposal does not close the neighboring projects' outstanding integer
indexing or simultaneous-universe composition queries merely by solving easier
numeric fold examples~\cite{leant,djex}.

\subsection{Certificate-guided, but not heuristic, semantic pruning}

Once universal equivalence is available for a restricted grammar, it can safely
remove redundant candidate implementations \emph{under an explicitly proved
observational congruence}. Suppose two complete candidates have the same numeric
output on every admissible input. They are interchangeable for a final
specification that observes only that output. They are not automatically
interchangeable as arbitrary higher-order subterms, because an enclosing context
might inspect their representation, intermediate state, dictionaries, or a
second output.

A safe first optimization deduplicates only \emph{complete candidates at the
root behavioral specification}, keeping one representative and its own source
receipt. A more ambitious grammar-level quotient needs a congruence theorem
for every grammar constructor. It must also preserve a representative for each
required source type and dictionary instantiation. The proposed workers supply
equivalence evidence; they do not make an arbitrary lossy hashing scheme safe.
This optimization is not part of the executed prototype.

\subsection{The first Lean implementation should emit small proofs}

There are two sensible replay paths. The initial path emits ordinary Lean
proofs of the finitely many certificate identities, using existing arithmetic
normalization and explicit substitution lemmas. The later path uses a reflected
sparse-polynomial checker with a proved evaluation theorem. Reflection is a
performance and proof-size improvement, not a prerequisite to validate the
first feature.

For Worker A, emit proofs of $q=tG$, $GA_a=C_aG$, and $Gs_0=0$. Apply a single
generic invariant theorem. For Worker B, prove the polynomial identities in
\eqref{eq:ideal-target}--\eqref{eq:ideal-initial}, evaluate them at a state,
and apply simultaneous invariant induction. The mathematical derivation of a
Gr\"obner basis need not appear in the Lean proof. For Worker C, prove the
local polynomial identity, apply the two binomial bridge lemmas, telescope a
finite sum, and only then address recurrence comparison and exceptional seeds.

The supplied \code{lean/ForgeExtension/Core.lean} is the semantic
list-execution induction glue. \code{Worked.lean} contains a hand-written replay
target for the mutually supporting nonlinear invariants and the local
squared-binomial identity. They contain no \code{sorry}, declared axiom, or
native-evaluation shortcut, but they were \emph{not compiled}. Their purpose is
to pin down the intended theorem shapes, not to pretend that the full reifier,
reflection theorem, or tactic exists.

\subsection{Concrete libraries and module responsibilities}

\begin{longtable}{@{}>{\raggedright\arraybackslash}p{0.23\linewidth}>{\raggedright\arraybackslash}p{0.35\linewidth}>{\raggedright\arraybackslash}p{0.34\linewidth}@{}}
\toprule
Component & Concrete choice & Required boundary\\
\midrule\endhead
Linear search & SymPy \code{Matrix.rank}, \code{gauss\_jordan\_solve}; later an
incremental echelon implementation & Exact rational coefficients only; no numeric-rank acceptance.\\
Ideal search & SymPy \code{groebner} over \code{QQ}, graded reverse lexicographic
order, \code{GroebnerBasis.reduce} & Export multiplier identities; the checker does
not trust basis claims or normal-form status.\\
Binomial synthesis & SymPy polynomial reduction and \code{Matrix.nullspace} &
Polynomial coefficient domain $\Q[A,B,n,k]$; retain endpoint-safe form.\\
Independent replay & Python standard library \code{fractions}, \code{json},
explicit sparse operations & No import of search code or symbolic library.\\
Lean polynomial semantics &
\path{Mathlib.Algebra.MvPolynomial.Eval} & Evaluation, substitution, and ring
homomorphism lemmas connect identities to source expressions.\\
Lean matrix semantics & \path{Mathlib.Data.Matrix.Mul} &
Matrix, vector, and dot-product conventions must agree with the serialized row/column layout.\\
Lean binomial semantics & \path{Mathlib.Data.Nat.Choose.Basic} &
Prove the zero-extended predecessor convention and rational casts explicitly.\\
General telescoping, optional & Sage \code{ore\_algebra} as an external proposer &
Only source-grounded pointwise identities and checked boundaries have proof authority.\\
\bottomrule
\end{longtable}

The SymPy API choices used by the prototypes are documented, and its polynomial
reduction returns quotient witnesses as well as a remainder~\cite{sympy-polys}.
The named Mathlib source files were inspected at Forge's pinned revision;
\code{MvPolynomial} provides evaluation maps, and the matrix file explicitly
specifies its row/vector conventions~\cite{mathlib-eval,mathlib-matrix}.
These are concrete integration targets, not claims that a compatible binary
build was available in the preparation environment.

The source package pins the same reference Lean and Mathlib revisions as Forge
for the uncompiled specimen project. A first real integration run must compile
against that pin, inspect all resulting axioms under Forge's existing policy,
and record the exact environment. Reusing Forge's policy is preferable to
inventing another almost-identical trust profile in this extension.
